% Generated draft: requires mathematical and source review.
\documentclass[11pt]{article}
\usepackage[margin=25mm]{geometry}
\usepackage{fontspec}
\setmainfont{Latin Modern Roman}
\usepackage{amsmath,amssymb,mathtools,unicode-math}
\setmathfont{Latin Modern Math}
\usepackage{enumitem,needspace,xurl,hyperref}
\setlength{\parindent}{0pt}
\setlength{\parskip}{5pt}
\setlength{\emergencystretch}{3em}
\newcommand{\N}{\mathbb{N}}
\newcommand{\Z}{\mathbb{Z}}
\newcommand{\Q}{\mathbb{Q}}
\newcommand{\R}{\mathbb{R}}
\newcommand{\C}{\mathbb{C}}
\newcommand{\F}{\mathbb{F}}
\newcommand{\A}{\mathbb{A}}
\newcommand{\PP}{\mathbb{P}}
\newcommand{\E}{\mathbb{E}}
\newcommand{\eps}{\varepsilon}
\newcommand{\dd}{\,\mathrm{d}}
\begin{document}
\section*{UnsolvedMath ID: 10900055}
\textbf{Source problem code:} catalogue entry
\subsection{Definitions and mathematical statement}
In the usual compact orientable setting, a Haken three-manifold is irreducible and contains a properly embedded two-sided incompressible surface of the required essential type (in the closed case, a closed positive-genus surface). A Heegaard splitting is \(M=V\cup_S W\), with \(V,W\) handlebodies in the closed case and the appropriate compression bodies when boundary is present. A compressing disk for \(S\) on one side is a properly embedded disk whose boundary is an essential simple closed curve in \(S\).

A splitting is strongly irreducible when both sides admit compressing disks and every pair of compressing disks on opposite sides intersects along its boundary:
\[\forall D\subset V\ \forall E\subset W,\qquad \partial D\cap\partial E\ne\varnothing.\]
Thus no disjoint opposite-side compressing-disk pair is allowed. Agol asks whether
\[\forall M\text{ Haken and hyperbolic},\quad
\exists\text{ a strongly irreducible Heegaard splitting of }M.\]
The quantifier is existence of at least one splitting for each manifold, not strong irreducibility of every splitting or of a specified minimum-genus splitting.

The source does not explicitly state closedness, finite volume or use of a compact core for a cusped manifold; in the boundary setting the Heegaard boundary partition/compression-body conventions must be fixed. Hyperbolicity in its standard complete three-manifold sense implies irreducibility; the research summary's parenthesis suggesting that Haken hyperbolic manifolds may be reducible is not a valid hypothesis to carry into this conventional question. The presence of incompressible surfaces is part of the Haken setting and is distinct from reducibility.
\subsection{Short English statement}
Does every Haken hyperbolic three-manifold admit at least one Heegaard splitting with no disjoint pair of compressing disks on opposite sides? Boundary and compactness conventions are not stated.
\subsection{Sources}
\begin{itemize}
\item Local UnsolvedMath registry, permanent ID \(10900055\), exact source code \texttt{catalogue entry}.
\item Agol, attribution in Delp, Hoffoss and Manning, \emph{Problems In Groups, Geometry, and Three-Manifolds} (2015), subsection 8.1, PDF page 6: \url{https://arxiv.org/pdf/1512.04620#page=6}; abstract \url{https://arxiv.org/abs/1512.04620}. Registry access 2026-07-29; source-tex extraction recorded as checked against rendered PDF text.
\item The research summary mentions Casson--Gordon and Rubinstein--Scharlemann theory and related non-Haken results, without exact citations or an exact theorem resolving this Haken existence question.
\end{itemize}
\end{document}
